\documentclass[aspectratio=169,12pt]{beamer}

% ---------- language & fonts ----------
\usepackage[utf8]{inputenc}
\usepackage[T1]{fontenc}
\usepackage[brazilian]{babel}
\usepackage{lmodern}
\renewcommand{\familydefault}{\sfdefault} % sans-serif throughout (accessibility)
\usepackage{fontawesome5}
\usepackage{tikz}
\usetikzlibrary{shapes.geometric,positioning,fit,arrows.meta}
\usepackage{pgfplots}
\pgfplotsset{compat=1.18}
\usepackage{array}
\usepackage{booktabs}
\usepackage{colortbl}
\usepackage{geometry}
% Beamer's native page is only ~16x9cm; PowerPoint's 16:9 slide is physically
% 13.333x7.5in. Matching that size makes literal pt sizes (32pt titles, 24pt
% body, per the accessibility spec) mean the same thing they mean in a deck
% authored in PowerPoint, instead of overflowing a much smaller native page.
\geometry{paperwidth=13.333in,paperheight=7.5in}

% ---------- palette (high contrast, colorblind-safe: navy / blue family / warm amber) ----------
\definecolor{navy}{HTML}{1E2761}
\definecolor{steel}{HTML}{4C6FB1}
\definecolor{ice}{HTML}{CADCFC}
\definecolor{amber}{HTML}{E8A33D}
\definecolor{muted}{HTML}{3E4A6B}

% ---------- beamer setup ----------
\usetheme{default}
\beamertemplatenavigationsymbolsempty
\setbeamercolor{background canvas}{bg=white}
\setbeamercolor{normal text}{fg=navy,bg=white}
\setbeamercolor{frametitle}{fg=navy,bg=white}
\setbeamercolor{itemize item}{fg=amber}
\setbeamercolor{itemize subitem}{fg=steel}
\setbeamertemplate{itemize item}{\textbullet}
\setbeamertemplate{itemize subitem}{\textendash}
\setbeamerfont{itemize item}{size=\normalsize}

% exact point sizes required by the accessibility guidelines: 32pt titles, 24pt body
\setbeamerfont{frametitle}{size={\fontsize{32}{37}\selectfont},series=\bfseries}
\setbeamerfont{normal text}{size={\fontsize{24}{29}\selectfont}}
\setbeamerfont{itemize/enumerate body}{size={\fontsize{24}{29}\selectfont}}
\setbeamerfont{itemize/enumerate subbody}{size={\fontsize{20}{25}\selectfont}}
\setbeamertemplate{frametitle}{%
  \vspace{4pt}\hspace{6pt}\usebeamerfont{frametitle}\usebeamercolor[fg]{frametitle}\insertframetitle\par
}
\setbeamertemplate{navigation symbols}{}
\setbeamersize{text margin left=8mm,text margin right=8mm}

% ---------- footer (page number + running title), always on white ----------
\newcommand{\decknote}{IHC 2026 \textperiodcentered\ Quixad\'a/CE --- Intera\c{c}\~ao Humano-Dados em Contextos Urbanos}
\setbeamertemplate{footline}{%
  \ifnum\insertframenumber=1\relax\else
    \leavevmode
    \hbox{\begin{beamercolorbox}[wd=\paperwidth,ht=3.4ex,dp=1.4ex,leftskip=6mm,rightskip=6mm]{footline}
      {\fontsize{9}{11}\selectfont\color{muted}\decknote\hfill\insertframenumber}
    \end{beamercolorbox}}
  \fi
}

% ---------- helper macros ----------
% NOTE on units: the deck's physical page is 13.333in x 7.5in (see \geometry
% above, matching a PowerPoint 16:9 slide), so all macro dimensions below are
% given in inches to be legible at that real physical scale.

% icon in a filled circle. #1=diameter (a full dimension, e.g. 1.2in) #2=bgcolor #3=icon command
\newcommand{\iconcircle}[3]{%
  \begin{tikzpicture}
    \node[circle,fill=#2,minimum size=#1,inner sep=0] (c) {};
    \node[white] at (c) {#3};
  \end{tikzpicture}%
}

% top-right decorative icon badge, placed via overlay so it doesn't disturb layout
\newcommand{\cornericon}[2]{% #1=bgcolor #2=icon
  \begin{tikzpicture}[remember picture,overlay]
    \node[anchor=north east,xshift=-9mm,yshift=-8mm] at (current page.north east) {\iconcircle{1in}{#1}{\LARGE #2}};
  \end{tikzpicture}%
}

% a short italic "tag" line under a title, e.g. the QP reference
\newcommand{\qptag}[1]{{\fontsize{16}{19}\selectfont\itshape\color{muted}#1\par}\vspace{2mm}}

% big-number stat callout: number on top, caption below, never overlapping
% (uses relative/chained positioning, like \pillarcard)
\newcommand{\statbox}[5]{% #1=width(in) #2=height(in) #3=bgcolor #4=number #5=caption
  \begin{tikzpicture}
    \node[rounded corners=4pt,fill=#3,minimum width=#1in,minimum height=#2in,text width=#1in,align=center] (b) {};
    \node[anchor=north,yshift=-8mm,navy,font=\bfseries\fontsize{40}{44}\selectfont] (n) at (b.north) {#4};
    \node[anchor=north,yshift=-3mm,navy,text width=#1in,align=center,font=\fontsize{16}{19}\selectfont,inner xsep=4mm] at (n.south) {#5};
  \end{tikzpicture}%
}

% pillar-style card: title + short text, on a navy/steel block, with icon
% uses relative (chained) positioning so it never overflows its own box
\newcommand{\pillarcard}[4]{% #1=bgcolor #2=icon #3=title #4=body
  \begin{tikzpicture}
    \node[rounded corners=4pt,fill=#1,minimum width=3.7in,minimum height=3.3in,text width=3.4in,align=center] (b) {};
    \node[anchor=north,yshift=-8mm] (ic) at (b.north) {\iconcircle{1.05in}{amber}{\Large #2}};
    \node[anchor=north,yshift=-4mm,text width=3.35in,align=center,white,font=\bfseries\fontsize{22}{25}\selectfont] (ti) at (ic.south) {#3};
    \node[anchor=north,yshift=-3mm,text width=3.1in,align=center,white,font=\fontsize{16}{19}\selectfont] at (ti.south) {#4};
  \end{tikzpicture}%
}

% auto-height highlight panel (safer than fixed-size tikz nodes for variable-length text)
% #1=width(w/ unit, e.g. 10.5in) #2=bgcolor #3=textcolor #4=title #5=body
\newcommand{\panel}[5]{%
  \colorbox{#2}{\begin{minipage}[t]{#1}
    \vspace{3mm}
    {\color{#3}\bfseries\fontsize{19}{23}\selectfont #4\par}
    \vspace{1.5mm}
    {\color{#3}\fontsize{15}{19}\selectfont #5\par}
    \vspace{3mm}
  \end{minipage}}%
}

\title{}
\author{}
\date{}

\begin{document}

% ============================================================ SLIDE 1 -- Titulo
{\setbeamercolor{background canvas}{bg=navy}
\begin{frame}[plain]
  \color{white}
  \begin{tikzpicture}[remember picture,overlay]
    \node[anchor=north east,xshift=-14mm,yshift=-12mm] at (current page.north east) {\iconcircle{1.5in}{amber}{\huge\faIcon{city}}};
  \end{tikzpicture}
  \vspace{18mm}

  {\fontsize{30}{34}\selectfont\bfseries\color{white}
   Intera\c{c}\~ao Humano-Dados\\ em Contextos Urbanos\par}
  \vspace{4mm}
  {\fontsize{20}{24}\selectfont\itshape\color{ice} Uma Revis\~ao Sistem\'atica da Literatura\par}
  \vspace{10mm}
  {\fontsize{18}{22}\selectfont\bfseries\color{white} La\'isa Dias Brito Alves \quad\textperiodcentered\quad Jo\~ao Luiz Bernardes Jr.\par}
  \vspace{2mm}
  {\fontsize{14}{18}\selectfont\color{ice} EACH-USP --- Escola de Artes, Ci\^encias e Humanidades \textperiodcentered\ Universidade de S\~ao Paulo\par}
  \vspace{14mm}
  {\fontsize{12}{15}\selectfont\color{ice} XXV Simp\'osio Brasileiro sobre Fatores Humanos em Sistemas Computacionais (IHC 2026) \textperiodcentered\ Quixad\'a/CE\par}
\end{frame}}

% ============================================================ SLIDE 2 -- Motivacao
\begin{frame}
  \frametitle{Cidades cada vez mais datificadas}
  \cornericon{amber}{\faIcon{city}}
  \vspace{3mm}
  \begin{columns}[T]
    \begin{column}{0.30\textwidth}
      \statbox{2.9}{2.3}{ice}{60\%}{da popula\c{c}\~ao mundial viver\'a em cidades at\'e 2050 (ONU)}
    \end{column}
    \begin{column}{0.66\textwidth}
      \begin{itemize}
        \item IoT e Urbanismo de Plataforma coletam dados em escala, com \textit{smartphones} como sensores passivos.
        \item Consentimento presumido (sem transpar\^encia real nem controle efetivo pelo cidad\~ao).
        \item Resultado: desafios cr\'iticos de privacidade e seguran\c{c}a na infraestrutura urbana.
      \end{itemize}
    \end{column}
  \end{columns}
\end{frame}

% ============================================================ SLIDE 3 -- O que e IHD
\begin{frame}
  \frametitle{Intera\c{c}\~ao Humano-Dados (IHD)}
  {\fontsize{20}{25}\selectfont Proposta por \textbf{Mortier et al. (2015)} --- campo aut\^onomo, focado na rela\c{c}\~ao \'etica entre o indiv\'iduo e seus rastros de dados.\par}
  \vspace{6mm}
  \begin{columns}[T]
    \column{0.333\textwidth}\centering
      \pillarcard{navy}{\faIcon{eye}}{Legibilidade}{Compreender e construir sentido sobre os dados coletados (\textit{sensemaking}).}
    \column{0.333\textwidth}\centering
      \pillarcard{navy}{\faIcon{handshake}}{Ag\^encia}{Poder real de decis\~ao sobre os pr\'oprios rastros digitais.}
    \column{0.333\textwidth}\centering
      \pillarcard{navy}{\faIcon{exchange-alt}}{Negociabilidade}{Revisar e revogar permiss\~oes conforme o contexto muda.}
  \end{columns}
\end{frame}

% ============================================================ SLIDE 4 -- Objetivo e Protocolo (metodo comprimido)
\begin{frame}
  \frametitle{Objetivo e protocolo da RSL}
  \vspace{4mm}
  \begin{center}
  \begin{tikzpicture}
    \node[rounded corners=6pt,fill=ice,minimum width=10.6in,minimum height=1.3in,text width=9.8in,align=center,text=navy,font=\bfseries\fontsize{24}{29}\selectfont] {Mapear a evolu\c{c}\~ao da IHD em contextos urbanos (2014--2024)};
  \end{tikzpicture}
  \end{center}
  \vspace{10mm}

  \begin{center}
  \begin{tikzpicture}
    \node[rounded corners=5pt,fill=navy,minimum width=2.8in,minimum height=2.0in] (b1) at (0,0) {};
    \node[white,font=\bfseries\fontsize{40}{44}\selectfont] at (b1.center) [yshift=6mm] {268};
    \node[white,text width=2.5in,align=center,font=\fontsize{16}{19}\selectfont] at (b1.center) [yshift=-9mm] {registros iniciais};

    \node[rounded corners=5pt,fill=steel,minimum width=2.8in,minimum height=2.0in] (b2) at (3.6in,0) {};
    \node[white,font=\bfseries\fontsize{40}{44}\selectfont] at (b2.center) [yshift=6mm] {170};
    \node[white,text width=2.5in,align=center,font=\fontsize{16}{19}\selectfont] at (b2.center) [yshift=-9mm] {exclu\'idos (CI e CE)};

    \node[rounded corners=5pt,fill=navy,minimum width=2.8in,minimum height=2.0in] (b3) at (7.2in,0) {};
    \node[white,font=\bfseries\fontsize{40}{44}\selectfont] at (b3.center) [yshift=6mm] {98};
    \node[white,text width=2.5in,align=center,font=\fontsize{16}{19}\selectfont] at (b3.center) [yshift=-9mm] {inclu\'idos no corpus};

    \draw[-{Latex[length=4mm]},line width=1.6pt,muted] (1.65in,0) -- (2.15in,0);
    \draw[-{Latex[length=4mm]},line width=1.6pt,muted] (5.25in,0) -- (5.75in,0);
  \end{tikzpicture}
  \end{center}
  \vspace{9mm}
  {\fontsize{16}{20}\selectfont\itshape\color{muted}\centering
   Protocolo: Kitchenham e Brereton (2013) \textperiodcentered\ RSL + PRISMA 2020 \textperiodcentered\ sem filtro de idioma \textperiodcentered\ sem restri\c{c}\~ao pr\'evia ao dom\'inio urbano (achado, n\~ao crit\'erio de entrada)\par}
\end{frame}

% ============================================================ SLIDE 5 -- Corpus + grafico
\begin{frame}
  \frametitle{O corpus: 98 estudos (2014--2024)}
  \vspace{2mm}
  \begin{center}
  \begin{tikzpicture}
    \begin{axis}[
      ybar stacked, bar width=22pt,
      width=10.8in, height=4.7in,
      ymin=0, ymax=25,
      symbolic x coords={2014,2015,2016,2017,2018,2019,2020,2021,2022,2023,2024},
      xtick=data, x tick label style={font=\fontsize{14}{16}\selectfont,navy},
      ytick={0,5,10,15,20,25}, y tick label style={font=\fontsize{14}{16}\selectfont,navy},
      ylabel={\fontsize{15}{18}\selectfont N\'umero de artigos}, ylabel style={navy},
      axis lines*=left, axis line style={navy,line width=0.8pt},
      enlarge x limits=0.06,
      legend style={at={(0.5,-0.20)},anchor=north,legend columns=3,font=\fontsize{15}{18}\selectfont,draw=none},
      ymajorgrids, grid style={gray!25},
    ]
      \addplot[fill=navy] coordinates {(2014,1)(2015,1)(2016,4)(2017,4)(2018,2)(2019,2)(2020,15)(2021,14)(2022,8)(2023,14)(2024,16)};
      \addplot[fill=steel] coordinates {(2014,0)(2015,0)(2016,0)(2017,0)(2018,1)(2019,3)(2020,2)(2021,2)(2022,2)(2023,1)(2024,3)};
      \addplot[fill=amber] coordinates {(2014,0)(2015,0)(2016,0)(2017,0)(2018,0)(2019,1)(2020,0)(2021,0)(2022,0)(2023,1)(2024,1)};
      \legend{Scopus (81), Science Direct (14), Teses CAPES (3)}
    \end{axis}
  \end{tikzpicture}
  \end{center}
  \vspace{2mm}
  {\fontsize{16}{20}\selectfont\itshape\color{muted}\centering
  Pico em 2020--2021 (33 estudos) e novo pico em 2024 (20 estudos) \textperiodcentered\ Scopus predomina em todos os anos ($\approx$82,7\%)\par}
\end{frame}

% ============================================================ SLIDE 6 -- QP1/QP4 Antes
\begin{frame}
  \frametitle{Antes da IHD: paradigma tecnocr\'atico}
  \cornericon{steel}{\faIcon{file-contract}}
  \qptag{QP1, QP4 --- Evolu\c{c}\~ao e Paradigmas}
  \panel{10.6in}{ice}{navy}{``Aviso e Escolha'' (\textit{Notice and Choice})}{Consentimento bin\'ario sobre termos extensos; EULAs unilaterais e inegoci\'aveis (\textit{take-it-or-leave-it}); a responsabilidade legal \'e transferida ao usu\'ario.}
  \vspace{5mm}
  \begin{itemize}
    \item Regula\c{c}\~ao inicial (Diretiva 95/46/CE $\rightarrow$ GDPR) tratava o dado --- replic\'avel e relacional --- como um objeto f\'isico isolado.
    \item Resultado: ecossistemas fechados (\textit{walled gardens}), que monopolizam a cust\'odia dos dados e restringem a interoperabilidade.
    \item A IHC cl\'assica limitava-se \`a usabilidade de tela, sem tratar os riscos \'eticos do processamento invis\'ivel e onipresente.
  \end{itemize}
\end{frame}

% ============================================================ SLIDE 7 -- QP1/QP4 Depois
\begin{frame}
  \frametitle{Depois: o dado como objeto de 1\textsuperscript{a} classe}
  \cornericon{amber}{\faIcon{layer-group}}
  \qptag{QP1, QP4 --- Evolu\c{c}\~ao e Paradigmas}
  \panel{10.6in}{navy}{white}{A virada de 2015 (Mortier et al.)}{A IHD emerge como campo \'etico aut\^onomo: dados deixam de ser entradas passivas e tornam-se entidades comput\'aveis --- manipul\'aveis, transfer\'iveis e audit\'aveis --- material de \textit{design}, n\~ao subproduto do sistema.}
  \vspace{5mm}
  \begin{itemize}
    \item \textbf{Legibilidade}: da exibi\c{c}\~ao est\'atica ao \textit{data sensemaking} --- mas jarg\~oes t\'ecnicos ainda persistem nas implementa\c{c}\~oes reais.
    \item \textbf{Ag\^encia}: do consentimento bin\'ario \`a \textit{accountability} --- o cidad\~ao vira \textit{steward} de seus rastros, por\'em sem heur\'isticas validadas de apoio.
    \item \textbf{Negociabilidade}: da tela ao IoT urbano onipresente --- o controle segue fragmentado entre m\'ultiplas ferramentas.
  \end{itemize}
\end{frame}

% ============================================================ SLIDE 8 -- QP2/QP3 Lacunas estruturais
\begin{frame}
  \frametitle{Lacunas estruturais e sobrecarga cognitiva}
  \cornericon{steel}{\faIcon{brain}}
  \qptag{QP2, QP3 --- Lacunas e Descompasso Pr\'atico}
  \begin{itemize}
    \item Falta de heur\'isticas de \textit{design} validadas, sobretudo em geoprocessamento --- lacuna que motivou o pr\'oprio \textit{framework} ATLAS.
    \item Ambiguidade na defini\c{c}\~ao dos perfis de usu\'ario e aus\^encia de crit\'erios claros do que torna um dado aberto reutiliz\'avel.
    \item Sobrecarga cognitiva negligenciada: fadiga e ansiedade informacional pela perda de controle sobre os pr\'oprios registros.
    \item Identidade Autossoberana (SSI): paradigma promissor de gest\~ao descentralizada de credenciais, mas ainda complexo demais para o p\'ublico leigo.
  \end{itemize}
\end{frame}

% ============================================================ SLIDE 9 -- QP2/QP3 IA e opacidade
\begin{frame}
  \frametitle{IA, opacidade algor\'itmica e o ``problema perverso''}
  \cornericon{steel}{\faIcon{robot}}
  \qptag{QP2, QP3 --- Lacunas e Descompasso Pr\'atico}
  \begin{itemize}
    \item Algoritmos de aprendizado de m\'aquina sofrem \textit{covariate shift}, dificultando o \textit{sensemaking} humano sobre fluxos de dados em tempo real.
    \item Sigilo industrial e caixas-pretas corporativas comprometem o Direito \`a Explica\c{c}\~ao de decis\~oes automatizadas que afetam o cidad\~ao.
    \item Servi\c{c}os de Gest\~ao de Dados Pessoais (PDMS) ainda falham em conciliar o extrativismo comercial com a real ag\^encia do usu\'ario.
  \end{itemize}
  \vspace{4mm}
  \panel{10.6in}{amber}{navy}{Um ``problema perverso''}{O equil\'ibrio entre otimizar servi\c{c}os em Cidades Inteligentes e proteger a privacidade \'e interdependente, sist\^emico e sem solu\c{c}\~ao puramente t\'ecnica.}
\end{frame}

% ============================================================ SLIDE 10 -- QP5/QPS4 Cidade civica
\begin{frame}
  \frametitle{Cidades Inteligentes e capacidade c\'ivica}
  \cornericon{amber}{\faIcon{broadcast-tower}}
  \qptag{QP5, QPS4 --- Cidades Inteligentes e Fatores Culturais}
  \begin{itemize}
    \item Tecnologia calma (\textit{calm technology}): monitorar o ambiente sem intrusão, com a privacidade como premissa central da arquitetura.
    \item Legibilidade c\'ivica: entender a coleta feita por sensores de rua liga atos cotidianos (como o transporte p\'ublico) \`as macroestruturas de extra\c{c}\~ao de dados.
    \item Ag\^encia via Limita\c{c}\~ao da Finalidade: direito de desativar a captura ambiental e de aprovar usos secund\'arios da informa\c{c}\~ao.
    \item \textit{Frameworks} locais --- HDIdea e ATLAS --- adaptam heur\'isticas e metáforas visuais ao contexto urbano espec\'ifico.
  \end{itemize}
\end{frame}

% ============================================================ SLIDE 11 -- QP5/QPS4 Fatores culturais e Brasil
\begin{frame}
  \frametitle{Fatores culturais e o contexto brasileiro}
  \cornericon{amber}{\faIcon{globe-americas}}
  \qptag{QP5, QPS4 --- Cidades Inteligentes e Fatores Culturais}
  \begin{itemize}
    \item Coletivismo $\times$ individualismo: culturas coletivistas aceitam mais o compartilhamento voltado ao benef\'icio comum; individualistas priorizam a restri\c{c}\~ao pessoal.
    \item Vi\'es de legado (\textit{legacy bias}): a importa\c{c}\~ao acr\'itica de modelos do Norte Global pode acentuar assimetrias no Brasil, marcado por disparidades de letramento digital.
    \item Minorias \'etnicas negociam menos seus dados com o Estado --- risco acrescido de discrimina\c{c}\~ao algor\'itmica e exclus\~ao das din\^amicas das Cidades Inteligentes.
    \item GranDIHC-BR (2025--2035): Grandes Desafios de Pesquisa nacionalizam a agenda da IHD no contexto brasileiro.
  \end{itemize}
\end{frame}

% ============================================================ SLIDE 12 -- QPS1-3 Multissensorial
\begin{frame}
  \frametitle{Do Mantra de Shneiderman ao multissensorial}
  \cornericon{steel}{\faIcon{hand-paper}}
  \qptag{QPS1--QPS3 --- Avalia\c{c}\~ao e Inclus\~ao Multissensorial}
  \begin{itemize}
    \item ``Mantra de Shneiderman'' (vis\~ao geral $\rightarrow$ \textit{zoom} $\rightarrow$ filtro) organiza pain\'eis, mas \'e insuficiente sozinho para o p\'ublico leigo.
    \item \textit{Scaffolding}: suporte progressivo ao aprendizado substitui sintaxes complexas por question\'amentos guiados.
    \item Interfaces de Voz (VUI) e sonifica\c{c}\~ao reduzem a barreira de aprendizado, explorando vias auditivas e vieses motores j\'a conhecidos.
    \item Fisicaliza\c{c}\~ao t\'atil e \textit{Databoxes} (nós de processamento local) tornam a governan\c{c}a de dados tang\'ivel.
  \end{itemize}
\end{frame}

% ============================================================ SLIDE 13 -- QPS1-3 Metodos de avaliacao
\begin{frame}
  \frametitle{M\'etodos de avalia\c{c}\~ao em IHD}
  \cornericon{steel}{\faIcon{clipboard-check}}
  \qptag{QPS1--QPS3 --- Avalia\c{c}\~ao e Inclus\~ao Multissensorial}
  \begin{itemize}
    \item \textit{Eye-tracking}: mede a sobrecarga cognitiva subconsciente em pain\'eis complexos (ex.: ATLAS).
    \item \textit{Think-Aloud}, entrevistas com peritos e o M\'etodo de Avalia\c{c}\~ao de Comunicabilidade (MAC) avaliam a compreens\~ao c\'ivica de forma pragm\'atica.
    \item Testes A/B, mapeamentos sistem\'aticos e pain\'eis \textit{Delphi} sustentam o consenso metodol\'ogico cont\'inuo da \'area.
    \item A comunidade brasileira de IHC nacionaliza esses protocolos como ferramenta contra as disparidades da dataficação.
  \end{itemize}
\end{frame}

% ============================================================ SLIDE 14 -- Matriz Comparativa
\begin{frame}
  \frametitle{Matriz comparativa de \textit{frameworks}}
  \vspace{2mm}
  \begin{center}
  {\fontsize{14}{17}\selectfont
  \renewcommand{\arraystretch}{1.35}
  \begin{tabular}{p{2.55in} p{2.15in} p{2.85in}}
    \toprule
    \textbf{\textit{Framework}} & \textbf{Princ\'ipio coberto} & \textbf{M\'etodo de avalia\c{c}\~ao} \\ \midrule
    Modelo Fundacional (Mortier et al.) & Legibilidade, Ag\^encia, Negociabilidade & Testes cl\'inicos, question\'arios, estudos de caso \\
    HDIdea & Legibilidade & \textit{Think-Aloud}, heur\'isticas colaborativas \\
    ATLAS (Geoprocessamento) & Legibilidade e Ag\^encia & \textit{Eye-tracking}, MAC \\
    \textit{Databoxes} / Mercados de Dados & Ag\^encia e Negociabilidade & Testes A/B, prototipa\c{c}\~ao tang\'ivel \\
    IHD Multissensorial (VUI e t\'atil) & Legibilidade (Acessibilidade) & Vieses motores, pain\'eis \textit{Delphi} \\
    \bottomrule
  \end{tabular}}
  \end{center}
  \vspace{4mm}
  {\fontsize{15}{19}\selectfont\itshape\color{muted}\centering
   Vers\~ao completa na Tabela 5 do artigo --- inclui tamb\'em \textit{Scaffolding} e Modelos H\'ibridos (CPS/XAI)\par}
\end{frame}

% ============================================================ SLIDE 15 -- Taxonomia por pilar
\begin{frame}
  \frametitle{Taxonomia de \textit{frameworks} por pilar}
  \vspace{2mm}
  \begin{columns}[T]
    \column{0.333\textwidth}\centering
      \pillarcard{navy}{\faIcon{eye}}{Legibilidade}{HDIdea, ATLAS, IHD Multissensorial, Diretrizes Cognitivas (\textit{Scaffolding}).}
    \column{0.333\textwidth}\centering
      \pillarcard{navy}{\faIcon{handshake}}{Ag\^encia}{ATLAS, \textit{Databoxes} e Mercados de Dados, Modelos H\'ibridos (CPS/XAI).}
    \column{0.333\textwidth}\centering
      \pillarcard{navy}{\faIcon{exchange-alt}}{Negociabilidade}{\textit{Databoxes} e Mercados de Dados, Modelos H\'ibridos (CPS/XAI).}
  \end{columns}
  \vspace{4mm}
  {\fontsize{15}{19}\selectfont\itshape\color{muted}\centering
   Contribui\c{c}\~ao original desta revis\~ao: ATLAS e os Modelos H\'ibridos cobrem mais de um pilar por sobreporem princ\'ipios\par}
\end{frame}

% ============================================================ SLIDE 16 -- Contribuicoes
\begin{frame}
  \frametitle{Contribui\c{c}\~oes desta revis\~ao}
  \cornericon{amber}{\faIcon{lightbulb}}
  \begin{itemize}
    \item Panorama atualizado (2014--2024) da IHD aplicada a contextos urbanos, com 98 estudos analisados sob o protocolo PRISMA 2020.
    \item Taxonomia in\'edita de \textit{frameworks} organizada pelos tr\^es pilares da IHD (Legibilidade, Ag\^encia, Negociabilidade).
    \item Matriz comparativa cruzando \textit{frameworks}, valida\c{c}\~ao emp\'irica e m\'etodos de avalia\c{c}\~ao.
    \item Mapeamento das Quest\~oes de Pesquisa aos eixos tem\'aticos e \`as lacunas sociot\'ecnicas identificadas.
  \end{itemize}
\end{frame}

% ============================================================ SLIDE 17 -- Limitacoes
\begin{frame}
  \frametitle{Limita\c{c}\~oes e trabalhos futuros}
  \cornericon{steel}{\faIcon{exclamation-triangle}}
  \begin{itemize}
    \item Triagem e extra\c{c}\~ao conduzidas pelos dois autores, sem protocolo formal de dupla checagem cega e independente.
    \item A \textit{string} de busca centrada em ``Human-Data Interaction'' pode ter omitido sin\^onimos (ex.: ``soberania digital'').
    \item Recorte temporal (2014--2024) e predom\'inio da Scopus ($\approx$82,7\%) limitam a generaliza\c{c}\~ao dos achados.
    \item \textbf{Trabalhos futuros}: valida\c{c}\~ao emp\'irica da taxonomia proposta; amplia\c{c}\~ao da \textit{string} de busca; estudos com grupos vulner\'aveis no contexto brasileiro.
  \end{itemize}
\end{frame}

% ============================================================ SLIDE 18 -- Obrigada
{\setbeamercolor{background canvas}{bg=navy}
\begin{frame}[plain]
  \color{white}
  \begin{tikzpicture}[remember picture,overlay]
    \node[anchor=north east,xshift=-14mm,yshift=-12mm] at (current page.north east) {\iconcircle{1.5in}{amber}{\huge\faIcon{heart}}};
  \end{tikzpicture}
  \vspace{28mm}
  {\fontsize{34}{38}\selectfont\bfseries\color{white} Obrigada!\par}
  \vspace{6mm}
  {\fontsize{20}{25}\selectfont\color{ice} Perguntas e coment\'arios s\~ao bem-vindos.\par}
  \vspace{14mm}
  {\fontsize{16}{20}\selectfont\bfseries\color{white} La\'isa Dias Brito Alves \quad\textperiodcentered\quad Jo\~ao Luiz Bernardes Jr.\par}
  \vspace{2mm}
  {\fontsize{13}{17}\selectfont\color{ice} EACH-USP --- Escola de Artes, Ci\^encias e Humanidades \textperiodcentered\ Universidade de S\~ao Paulo\par}
\end{frame}}

\end{document}
